\documentclass[11pt]{article}

\usepackage[margin=1in]{geometry}
\usepackage[T1]{fontenc}
\usepackage{lmodern}
\usepackage{microtype}
\usepackage{amsmath,amssymb}
\usepackage{graphicx}
\usepackage{booktabs}
\usepackage{array}
\usepackage{multirow}
\usepackage{xcolor}
\usepackage{enumitem}
\usepackage[hidelinks]{hyperref}

% Powered real-LLM run parameters (6 collapse-capable tasks x 8 repeats).
\newcommand{\poweredTasks}{6}
\newcommand{\poweredRepeats}{8}
\newcommand{\poweredEpisodes}{576}

\title{\textbf{When the Governor Becomes the Disturbance:}\\[0.3em]
Control-Generated Disturbance and Cost-Aware Backoff in Governed Tool-Using Agents}
\author{Veronique Ziegler\\Independent Researcher, Critical Attention Systems}
\date{\today}

\begin{document}
\maketitle

\begin{abstract}
A governor that monitors a tool-using agent and intervenes to keep it on task
is normally treated as a pure stabiliser. We show it need not be. When
corrective intervention carries a cost---when tightening the control loop
itself perturbs the actions it regulates---a cost-blind governor can drive the
agent it supervises into a self-reinforcing collapse: each induced failure
raises the regulatory state, which tightens control further, which induces more
failures. We study this \emph{control-generated disturbance} in a discrete
tool-use substrate. On a large-sample stochastic-policy substrate we
demonstrate three things: (i) raising the induced-cost scale $\alpha$ under a
cost-blind governor amplifies into a block-storm collapse, and this holds under
two independent cost channels---a fabricated transient failure and a
genuinely-executed corrupted action---so the collapse is not an artifact of the
injection; (ii) a \emph{cost-aware} governor that backs off when its own recent
tightening has been generating failures suppresses its regulatory intensity in
a monotone dose-response with backoff sensitivity; and (iii) against a full
baseline ladder, cost-aware backoff beats a fixed weak governor matched to its
realised intervention rate, so the benefit is genuinely adaptive rather than
merely intervening less. On a real large-language-model agent
(Gemini~2.5~Flash), a powered run ($\poweredTasks$ collapse-capable tasks
$\times$ $\poweredRepeats$ repeats, $\poweredEpisodes$ episodes) reproduces all
three effects: amplification into block-storm collapse, the suppression
dose-response, and success rescue that peaks at an intermediate backoff
strength. The results identify intervention cost
and cost-awareness as design variables for agent governors: a governor's own
interventions, when costly, can create positive feedback that a cost-aware
governor breaks---a structure that appears on both a hand-coded substrate and a
real LLM, so it is not an artifact of any single formalism.
\end{abstract}

\section{Introduction}

Autonomous tool-using agents are increasingly deployed with a supervisory
\emph{governor}~\cite{astrom2008feedback}: a monitor that inspects the agent's
trajectory, emits a regulatory signal, and can warn or block proposed actions
to keep the agent
from looping, retrying failed assumptions, or answering prematurely. The
implicit assumption is that governance is one-directionally good: more
regulation can only help, or at worst do nothing.

This assumption fails when regulatory intervention has a cost that feeds back
into the agent's own state. A governor that tightens its control in response to
trouble is, in realistic systems, not free: a blocked action forces a re-plan,
an injected constraint narrows the action distribution, an aggressive
intervention can corrupt or delay the next tool call. If that induced
disturbance then registers as further trouble, the governor tightens again, and
the loop can run away. The agent is no longer merely being regulated;
\emph{the governor has become a disturbance source}.

This question is a natural extension of the IRAM-$\Omega$-Q line of
adaptive-agent
models~\cite{ziegler2026iramomegaq,ziegler2026carryover,ziegler2026switching},
which studies, in a continuous quantum-like state model, how the timing and
demand of regulation shape an agent's control burden. It is also a standard
concern in feedback control, where an over-eager regulator can inject the very
disturbance it is trying to correct~\cite{astrom2008feedback}. We pose it as a
hypothesis---that a governor's own tightening, when it carries a cost, becomes a
disturbance source that can drive positive feedback---and ask whether it holds
on two unrelated substrates: a symbolic tool-using agent supervised by a
finite-state governor, and a real large-language-model (LLM) agent. If the
positive-feedback-and-cost-awareness structure appears under both, it is not an
artifact of any single formalism.

Our results are demonstrated on two substrates: a hand-coded stochastic-policy
substrate, where matched large-$N$ seeding is cheap, and a powered run on a real
LLM:
\begin{enumerate}[leftmargin=1.4em]
\item a cost-blind governor amplifies into block-storm collapse as $\alpha$
grows, on both substrates, under two independent induced-cost channels
(Sections~\ref{sec:amplification},~\ref{sec:gemini});
\item a cost-aware governor suppresses its intensity in a monotone
dose-response with backoff sensitivity, on both substrates
(Sections~\ref{sec:suppression},~\ref{sec:gemini}); and
\item against a full baseline ladder, cost-aware backoff outperforms a fixed
weak governor at matched intervention rate, establishing that the effect is
adaptive (Section~\ref{sec:baselines}).
\end{enumerate}
On the real LLM, the powered run (Section~\ref{sec:gemini}) confirms the full
collapse-and-rescue picture---amplification, suppression, and rescue---and adds
a design finding: the rescue peaks at an intermediate backoff strength.

\section{Background and scope}
\label{sec:background}

The IRAM-$\Omega$-Q model represents an adaptive agent by a normalised state,
an entropy-like uncertainty signal, and an adaptive gain $\mu(t)$ interpreted
as regulatory demand, and compares regulation-first (RF) and disturbance-first
(DF) causal orderings~\cite{ziegler2026carryover}. A recurring theme in that
line, and in feedback control more generally~\cite{astrom2008feedback}, is that
the \emph{source} of a disturbance matters, not only its magnitude---and in
particular that a controller's own tightening can itself act as a disturbance.
We model that tightening cost, in the continuous formalism's terms, as an
induced disturbance
$\eta_{\mathrm{induced}} =
\alpha\max(0,\mu_{\mathrm{after}}-\mu_{\mathrm{before}})^2$ charged only when
the controller tightens, and take it as the mechanism whose discrete analogue we
study.

Our hypothesis is exactly this---that control-generated disturbance can create
positive feedback that a cost-aware controller can break---tested on a different
substrate. We do \emph{not} claim to reproduce the continuous formalism's
numerics or architecture; the discrete induced-cost mechanism of
Section~\ref{sec:method} is a deliberately simple analogue of
$\eta_{\mathrm{induced}}$, and the correspondence is structural, not
quantitative.

\section{The governed tool-use substrate}
\label{sec:substrate}

\paragraph{Agent, governor, task.}
An episode places an agent in a small file workspace and asks it to locate and
summarise a target file; the prompt embeds an incorrect assumption (a wrong
filename, folder, extension, or a confidently-named nonexistent path), so
recovery requires inspecting the workspace. Each step the agent proposes one
tool action (\texttt{READ\_FILE} or \texttt{LIST\_DIRECTORY}) or a final
answer. A finite-state \emph{governor} evaluates the trajectory and emits a
regulatory state, ordered by severity,
$\textsc{Normal}\to\textsc{Disturbed}\to\textsc{Recovering}\to\textsc{Stuck}$,
plus a terminal \textsc{Released}. \textsc{Disturbed} follows any failed tool
call; \textsc{Recovering} and \textsc{Stuck} follow repeated failures or
sustained looping; \textsc{Released} follows recovery of sufficient content.
The governor can reshape the agent's action distribution toward inspection
(reweighting) and can \emph{block} an action that repeats a failed assumption
or continues tool use after \textsc{Released}. Task success is a substring
match of the final answer against the known target answer.

\paragraph{Two agent realisations.}
The \textbf{stochastic policy}~\cite{sutton2018rl} samples each step from a governor-modulated
weighted categorical over candidate actions; behavioural presets give
archetypes (an \emph{impulsive} high-entropy sampler; a \emph{persistent}
archetype that clings to failed actions). It is deterministic-in-distribution,
so thousands of matched seeds are cheap; it tests substrate-independence of the
mechanism, not learned behaviour. The \textbf{real LLM agent} delegates action
selection to Gemini~2.5~Flash via a JSON tool-call prompt embedding the current
regulatory signal.

\paragraph{Regulation intensity.}
We give the governor's tightening a scalar via a monotone intensity over its
state,
\begin{equation}
\begin{gathered}
g(\textsc{Normal})=0,\quad
g(\textsc{Disturbed})=0.75,\quad
g(\textsc{Recovering})=0.85,\\
g(\textsc{Stuck})=0.95,\quad
g(\textsc{Released})=0.99,
\end{gathered}
\label{eq:intensity}
\end{equation}
equal to one minus the governor's retry-suppression multiplier. Higher $g$
means harder clamping; this is the discrete counterpart of the continuous gain.

\section{Method: control-generated disturbance}
\label{sec:method}

\paragraph{Induced cost.}
Only \emph{positive} tightening generates cost. With
$\Delta g^{+}(t)=\max(0,\,g(t)-g(t-1))$, an induced event is drawn at each step
with probability
\begin{equation}
p_{\mathrm{induced}}(t)=\min\!\big(1,\;\alpha\,\Delta g^{+}(t)^2\big),
\label{eq:pinduced}
\end{equation}
$\alpha\ge0$ the single swept scale. At $\alpha=0$ no cost is ever generated
and the agent reduces exactly to the ordinary governed agent (a regression
anchor). Because the escalation $\textsc{Normal}\to\textsc{Disturbed}$ is a
large jump ($\Delta g^{+}=0.75$), firing saturates near $\alpha\approx1.8$, so
the informative range is $\alpha\in[0,2]$.

\paragraph{Two cost channels.}
An induced event is realised through one of two mechanisms, applied to the
agent's intended action so that the governor's repeat logic engages
identically:
\begin{itemize}[leftmargin=1.4em,itemsep=2pt]
\item \textbf{Result injection}: the intended action's result is replaced by a
fabricated transient failure.
\item \textbf{Argument corruption}: a bad-path (corruption-suffixed) version of
the intended action is \emph{actually executed} by the real tool, which returns
its own genuine not-found error; the failure is attributed to the intended
action. Nothing is fabricated---a genuinely degraded call really failed.
\end{itemize}
Reporting both guards against the objection that the collapse is engineered by
a single injection trick. (A softer corruption that leaves the target valid
attenuates the loop, confirming that the cost must yield governor-visible
failures to drive escalation.)

\paragraph{The feedback loop.}
An induced failure is a real failed action in history. The governor escalates,
tightens ($g$ rises), which raises $\Delta g^{+}$ and hence
$p_{\mathrm{induced}}$ at later escalations, producing more failures. A
cost-blind governor always escalates on failure and cannot exit the loop; the
collapse manifests as a \emph{block storm}, because the governor also blocks
the agent from retrying its now-failed intended action.

\paragraph{Cost-aware backoff.}
The cost-aware governor maintains an exponentially weighted induced-failure rate
$r(t)=(1-\lambda)r(t-1)+\lambda\,\mathbb{1}[\text{induced at }t]$
($\lambda=0.3$) and applies its regulation---reweighting \emph{and}
blocking---only with probability
\begin{equation}
p_{\mathrm{apply}}(t)=1-\min\!\big(1,\;k\,r(t)\big),
\label{eq:papply}
\end{equation}
$k>0$ the backoff sensitivity (swept over $\{1.5,3.0,6.0\}$). On a backed-off
step the agent sees \textsc{Normal} and no block fires, so $g=0$, no cost is
generated, $r$ decays, and $p_{\mathrm{apply}}$ recovers---negative feedback
that can break the loop. A \textbf{fixed weak governor} that applies regulation
at a constant probability, non-adaptively, serves as a control for whether
adaptivity itself matters (Section~\ref{sec:baselines}).

\section{Demonstrated results I: stochastic substrate}
\label{sec:stochastic}

All results in this section use $2000$ matched seeds per (task, condition)
cell across eight file-recovery tasks, real-time (RF) governance, and a
16-step budget.

\subsection{Cost-blind governance amplifies into collapse}
\label{sec:amplification}

Under the cost-blind governor, raising $\alpha$ collapses success while the
governor's intensity stays pinned, and this holds under both cost channels
(Table~\ref{tab:mechanisms}). For the persistent archetype, success falls
$0.88\to0.49\to0.18$ (result injection) and $0.88\to0.48\to0.19$ (argument
corruption) as $\alpha$ goes $0\to1\to2$, with $\bar g$ pinned near $0.71$
throughout. The near-identity across a fabricated-failure channel and a
genuinely-executed-corrupted-action channel shows the collapse is a property of
induced governor-visible failure, not of the specific injection.

\begin{table}[t]
\centering
\caption{Mechanism robustness (persistent archetype, $2000$ seeds/cell $\times$
$8$ tasks). Success and mean intensity $\bar g$ by cost channel, backoff, and
$\alpha$. Cost-blind amplification and cost-aware suppression are essentially
identical under both channels.}
\label{tab:mechanisms}
\begin{tabular}{llccc}
\toprule
 & & $\alpha{=}0$ & $\alpha{=}1$ & $\alpha{=}2$ \\
\midrule
\multirow{2}{*}{Result injection}
 & cost-blind (succ / $\bar g$)  & $0.88\,/\,0.72$ & $0.49\,/\,0.71$ & $0.18\,/\,0.70$ \\
 & cost-aware (succ / $\bar g$)  & $0.88\,/\,0.72$ & $0.65\,/\,0.50$ & $0.40\,/\,0.29$ \\
\midrule
\multirow{2}{*}{Argument corruption}
 & cost-blind (succ / $\bar g$)  & $0.88\,/\,0.72$ & $0.48\,/\,0.71$ & $0.19\,/\,0.70$ \\
 & cost-aware (succ / $\bar g$)  & $0.88\,/\,0.72$ & $0.64\,/\,0.50$ & $0.40\,/\,0.29$ \\
\bottomrule
\end{tabular}
\end{table}

\subsection{Cost-aware governance suppresses its own intensity}
\label{sec:suppression}

Under the cost-aware governor, $\bar g$ falls with $\alpha$ (Table~\ref{tab:mechanisms},
persistent: $0.72\to0.50\to0.29$) while the cost-blind $\bar g$ stays pinned---the
suppression branch, absent by construction under cost-blindness. Suppression
deepens monotonically with backoff sensitivity $k$ (impulsive $k{=}6$:
$0.36\to0.11$; persistent $k{=}6$: $0.69\to0.21$). Success rescue is
agent-dependent: the persistent archetype is rescued (Table~\ref{tab:mechanisms},
$\alpha{=}2$: cost-blind $0.18$ vs cost-aware $0.40$), whereas the impulsive
archetype---whose un-governed floor is near-catastrophic---is not, because
backing off returns it toward that floor faster than it saves on induced cost.
Cost-awareness helps exactly when the agent has a survivable un-regulated floor.

\subsection{Cost-aware backoff is adaptive, not merely less governance}
\label{sec:baselines}

Table~\ref{tab:baselines} places the induced-cost conditions on a full baseline
ladder at $\alpha=2$. The key control is the \emph{fixed weak governor}, which
intervenes at a constant probability matched to cost-aware's realised
intervention rate ($0.45$ for persistent). For the persistent archetype:
ordinary governance at $\alpha=0$ lifts success far above no governance
($0.88$ vs $0.31$); induced cost makes cost-blind governance actively harmful,
dropping success \emph{below} no governance ($0.19$); cost-aware backoff
restores it above the un-governed level ($0.40$); and, decisively, the fixed
weak governor at the same intervention rate reaches only $0.30$. Adaptivity---%
intervening when it matters and standing down when its tightening is
backfiring---accounts for the $0.10$ gap between cost-aware and rate-matched
fixed governance; it is not reducible to governing less. The impulsive
archetype shows the same ladder shape with the rescue absent, consistent with
its non-survivable un-governed floor.

\begin{table}[t]
\centering
\caption{Baseline ladder at $\alpha=2$ (result injection, $2000$ seeds/cell
$\times$ $8$ tasks). The fixed weak governor is matched to cost-aware's realised
intervention (apply) rate. For persistent, cost-aware ($0.40$) beats
rate-matched fixed ($0.30$), and cost-blind governance ($0.19$) is worse than no
governance ($0.31$).}
\label{tab:baselines}
\begin{tabular}{lcccc}
\toprule
 & \multicolumn{2}{c}{persistent} & \multicolumn{2}{c}{impulsive} \\
\cmidrule(lr){2-3}\cmidrule(lr){4-5}
condition & success & apply rate & success & apply rate \\
\midrule
no governor                 & $0.31$ & --     & $0.09$ & -- \\
governor, $\alpha{=}0$       & $0.88$ & $1.00$ & $0.51$ & $1.00$ \\
cost-blind, $\alpha{=}2$     & $0.19$ & $1.00$ & $0.11$ & $1.00$ \\
cost-aware, $\alpha{=}2$     & $\mathbf{0.40}$ & $0.46$ & $0.10$ & $0.75$ \\
fixed weak, $\alpha{=}2$ (matched) & $0.30$ & $0.45$ & $0.08$ & $0.72$ \\
\bottomrule
\end{tabular}
\end{table}

\section{Demonstrated results II: real-LLM amplification, suppression, and rescue}
\label{sec:gemini}

We ran the induced-cost sweep with Gemini~2.5~Flash as the agent: a cost-blind
governor and cost-aware governors at backoff sensitivity $k\in\{1.5,3.0,6.0\}$,
induced-cost scales $\alpha\in\{0,1,2\}$, six collapse-capable wrong-path tasks (each escalation-prone \emph{and} solvable---a
confident wrong-path hint paired with a short verbatim-reproducible answer), and
eight repeats per cell, for $\poweredEpisodes$ episodes with per-episode seeds
avalanche-finalised for independence (free-tier limits handled by checkpointed
resumption and backoff on transient errors). Pooled results exclude one task
whose answer never substring-matched; the remaining five engaged the governor
($\bar g>0.4$ at $\alpha=0$) with high baseline success, so the induced cost has
room both to collapse success and to be rescued. An initial three-repeat
$k$-sweep established suppression but was underpowered for collapse and
rescue---at $\alpha=1$ the induced draw is a ${\approx}0.56$ per-escalation coin
flip that three repeats cannot resolve---and it exposed a seeding defect that
had made repeats non-independent; the run reported here fixes both. All three
effects are present and well powered (Table~\ref{tab:gemini}).

\paragraph{Amplification.} Under the cost-blind governor, success collapses
monotonically with $\alpha$, $0.83\to0.58\to0.33$, while mean hard blocks rise
$0.1\to4.0\to7.6$ per episode and the answered rate falls to $0.35$: on the
three tasks that storm, blind reaches ${\sim}13$ hard blocks per episode and
most runs never produce an answer. The governor's own tightening buries the
agent.

\paragraph{Suppression.} Mean intensity $\bar g$ falls with $\alpha$ under the
cost-aware governor and stays pinned under the cost-blind one, in a clean
monotone dose-response with backoff sensitivity: the $\bar g$ drop from
$\alpha{=}0$ to $\alpha{=}2$ is $0.25,0.43,0.49$ for $k=1.5,3.0,6.0$
(Table~\ref{tab:gemini}).

\paragraph{Rescue, with an optimum in $k$.} At $\alpha=2$ the cost-blind
governor sits at $0.33$ success; every cost-aware mode beats it, and the rescue
peaks at an \emph{intermediate} backoff strength: $0.50,0.75,0.71$ for
$k=1.5,3.0,6.0$, so $k{=}3$ gives the largest rescue ($+0.42$ over blind). Too
gentle a backoff ($k{=}1.5$) leaves the block storm partly intact; too
aggressive ($k{=}6$) stands down governance the agent still needs. The rescue is
broad-based: of the three tasks that collapse under blind (T002, E003, C001;
$0$--$12\%$ success, ${\sim}12$ hard blocks each), all are restored by $k{=}3$
backoff to $88$--$100\%$ with block storms broken to near zero
(Table~\ref{tab:gemini-tasks}); the two tasks that never stormed under blind had
little to rescue. As on the stochastic substrate, the rescue mechanism is the
breaking of the block loop, not avoidance of induced cost---cost-aware runs
often incur \emph{more} induced failures because they keep acting.

\begin{table}[t]
\centering
\caption{Real-LLM results (Gemini, $\poweredEpisodes$-episode powered run,
$n=48$ per cell, pooled over five collapse-capable tasks). Cost-blind amplifies
(success falls, blocks storm, intensity pinned); cost-aware suppresses its
intensity in a monotone $k$ dose-response and rescues success, with the rescue
peaking at $k{=}3$.}
\label{tab:gemini}
\begin{tabular}{llcccc}
\toprule
 & $\alpha$ & blind & $k{=}1.5$ & $k{=}3.0$ & $k{=}6.0$ \\
\midrule
\multirow{3}{*}{success}
 & $0$ & $0.83$ & $0.83$ & $0.83$ & $0.83$ \\
 & $1$ & $0.58$ & $0.75$ & $0.77$ & $0.77$ \\
 & $2$ & $0.33$ & $0.50$ & $\mathbf{0.75}$ & $0.71$ \\
\midrule
\multirow{3}{*}{intensity $\bar g$}
 & $0$ & $0.66$ & $0.69$ & $0.69$ & $0.68$ \\
 & $1$ & $0.71$ & $0.62$ & $0.47$ & $0.40$ \\
 & $2$ & $0.74$ & $0.44$ & $0.26$ & $0.19$ \\
\midrule
hard blocks & $2$ & $7.6$ & $2.4$ & $0.3$ & $0.1$ \\
\bottomrule
\end{tabular}
\end{table}

\begin{table}[t]
\centering
\caption{Per-task success and hard blocks at $\alpha=2$ (Gemini, blind vs
$k{=}3$ cost-aware, $n=8$ per cell). The three tasks that storm under blind are
fully rescued, with block storms broken.}
\label{tab:gemini-tasks}
\begin{tabular}{lcccc}
\toprule
 & \multicolumn{2}{c}{success} & \multicolumn{2}{c}{hard blocks} \\
\cmidrule(lr){2-3}\cmidrule(lr){4-5}
task & blind & $k{=}3$ & blind & $k{=}3$ \\
\midrule
T002 & $0.00$ & $\mathbf{1.00}$ & $12.8$ & $0.2$ \\
E003 & $0.00$ & $\mathbf{0.88}$ & $12.8$ & $0.8$ \\
C001 & $0.12$ & $\mathbf{1.00}$ & $11.4$ & $0.4$ \\
E002 & $0.88$ & $1.00$ & $1.8$ & $0.0$ \\
C002 & $1.00$ & $0.62$ & $1.6$ & $0.1$ \\
\bottomrule
\end{tabular}
\end{table}

\section{Discussion}

The demonstrated results support one transferable claim: \emph{intervention cost
and cost-awareness are design variables for agent governors}. A governor whose
corrective actions perturb the agent is not a safe default; at sufficient
intervention cost a cost-blind governor can be worse than no governor at all
(Table~\ref{tab:baselines}), collapsing the agent it was meant to protect
through a self-reinforcing block storm. The remedy is not weaker governance but
\emph{cost-aware} governance that senses when its own tightening is backfiring
and stands down---and the fixed-weak control shows this benefit is adaptive, not
an artifact of intervening less. That the amplification reproduces under two
distinct cost channels, and that the suppression dose-response reproduces on a
real LLM, indicates the effect is not tied to the specific injection or to the
continuous IRAM-$\Omega$-Q formalism.

We are deliberately narrow about the connection to that formalism. The
defensible statement is that a single structural claim---control-generated
disturbance can create positive feedback, and cost-awareness can break
it---holds across two very different substrates; we do not claim to validate the
IRAM-$\Omega$-Q architecture or reproduce its numerics. Two practical nuances follow from
the data: cost-aware backoff helps only when the agent has a survivable
un-regulated floor, and the optimal backoff strength must be tuned---it is
intermediate and substrate-dependent, gentlest on the stochastic policy and
peaking at $k{\approx}3$ on the real LLM, where too little backoff leaves the
block storm intact and too much stands down governance the agent still needs.

\section{Limitations}

The stochastic-policy agent is hand-coded, so those (large-$N$) results bound
substrate-independence of the mechanism, not learned-agent behaviour. On the
real LLM, all three effects are powered ($n=48$ per cell), but the results come
from a single model (Gemini~2.5~Flash) and five collapse-capable tasks; a
broader task suite and additional models would strengthen generality. The
induced-cost model is a one-sided quadratic in positive
tightening; other cost shapes and injection points are possible, and a softer
argument corruption that leaves targets valid attenuates the loop. The intensity
ladder~\eqref{eq:intensity} is a monotone proxy for clamp strength, not a
calibrated cost, and correctness is a substring match, which concentrates the
success signal on tasks with short answers.

\section{Conclusion}

A supervisory governor for a tool-using agent can itself become a disturbance
source when its interventions carry a cost that feeds back into the agent's
state. On a large-sample stochastic substrate we demonstrate that a cost-blind
governor amplifies this into block-storm collapse---under two independent cost
channels, and to the point of being worse than no governor---while a cost-aware
governor that backs off suppresses its intensity in a monotone dose-response and
outperforms a rate-matched fixed governor, so the benefit is adaptive. A powered
run on a real LLM agent (Gemini~2.5~Flash) reproduces all three effects---%
amplification, the suppression dose-response, and rescue---with the rescue
peaking at an intermediate backoff strength. Governor intervention cost, and the
governor's awareness of it, are design variables that fixed-strength governance
overlooks.

\section*{Data and code availability}
The governed-agent framework, induced-cost agents (both cost channels and the
backoff regimes), experiment harnesses, and analysis scripts are maintained in
the project repository; numerical summaries supporting the reported tables will
be released with the powered confirmation run.

\begin{thebibliography}{9}

\bibitem{ziegler2026iramomegaq}
V.~Ziegler.
\newblock \emph{IRAM-$\Omega$-Q: A Computational Framework for Uncertainty
Regulation in Adaptive Agents}.
\newblock arXiv:2603.16020, 2026.

\bibitem{ziegler2026carryover}
V.~Ziegler.
\newblock \emph{Anticipatory Regulation, Control Demand, and Hysteresis in
Artificial Agents}.
\newblock arXiv:2606.30975, 2026.

\bibitem{ziegler2026switching}
V.~Ziegler.
\newblock \emph{Intermittent Control Is Not Diluted Control: A Switching Effect
in Artificial Agency}.
\newblock arXiv:2607.17432, 2026.

\bibitem{astrom2008feedback}
K.~J. \r{A}str\"om and R.~M. Murray.
\newblock \emph{Feedback Systems: An Introduction for Scientists and Engineers}.
\newblock Princeton University Press, 2008.

\bibitem{sutton2018rl}
R.~S. Sutton and A.~G. Barto.
\newblock \emph{Reinforcement Learning: An Introduction}.
\newblock MIT Press, second edition, 2018.

\end{thebibliography}

\end{document}
